\scanpage{30}
\mspage{\handwritten{8}}
comme ci-dessus la fonctorialité de $\handwritten{\alpha}$ en $P$ et en $Q$, on est donc ramené au cas où on a \deleted{$\unsure{P=\underline{\mathrm Z}/n}$}
\[
P=Q=(\underline{\mathrm Z}/n\underline{\mathrm Z})_{\underline{T}}\ ,
\]
$p$ et $q$ étant la section unité de $(\underline{\mathrm Z}/n\underline{\mathrm Z})_{\underline{T}}$, et on est donc ramené à déterminer l'homomorphisme
\[
\handwritten{\alpha}:(\underline{\mathrm Z}/n\underline{\mathrm Z})_{\underline{T}}\otimes(\underline{\mathrm Z}/n\underline{\mathrm Z})_{\underline{T}}\handwritten{\longrightarrow}\underline{\mathrm{Tor}}_1^{\underline{\mathrm Z}}((\underline{\mathrm Z}/n\underline{\mathrm Z})_{\underline{T}},(\underline{\mathrm Z}/n\underline{\mathrm Z})_T)\ .
\]
La description 2.1.7 nous montre que c'est l'isomorphisme
\[
(\underline{\mathrm Z}/n\underline{\mathrm Z})_{\handwritten{\underline T}}\handwritten{\simeq}\underline{\mathrm{Tor}}_1^{\underline{\mathrm Z}}((\underline{\mathrm Z}/n\underline{\mathrm Z})_{\underline{T}},(\underline{\unsure{\mathbb Z|\mathrm Z}}/n\underline{\mathrm Z})_T)\tag*{(2.1.8.1)}
\]
déduit de la résolution plate canonique
\[
0\handwritten{\longrightarrow}\underline{\mathrm Z}_T\mathrel{\mathop{\handwritten{\longrightarrow}}\limits^{n}}\underline{\mathrm Z}_T\handwritten{\longrightarrow}0\tag*{\handwritten{(2.1.8.2)}}
\]
% circled
\margin{\handwritten{$(\mathbb Z/n\mathbb Z)_{\underline T}$}}
du \uline{premier} facteur,\deleted{dans} \inserted{$P=(\underline{\mathrm Z}/n\underline{\mathrm Z})_{\handwritten{T}},$} qui figure dans le $\underline{\mathrm{Tor}}_1$ .

2.1.9. La définition donnée de (2.1.1) n'est pas symétrique, et on peut, en échangeant les rôles de $P$ et de $Q$ dans la construction précédente, définir un homomorphisme analogue
\[
\handwritten{\beta^n_{P,Q}}:\ {}_nP\otimes{}_nQ\handwritten{\longrightarrow}\underline{\mathrm{Tor}}_1(P,Q)\ ,\tag*{(2.1.9.1)}
\]
en se ramenant d'abord au cas $nP=0$, et utilisant une résolution plate \uline{de} $Q$. On peut aussi définir $\handwritten{\beta}$ en termes de $\handwritten{\alpha}$, de façon équivalente, par la commutativité du diagramme
\diagram{\begin{tikzcd}[ampersand replacement=\&,column sep=large,row sep=large]
\deleted{\text{(2}} \& {}_nP\otimes{}_nQ\handwritten{\arrow[r,"\sim"]}\handwritten{\arrow[d,"{\beta^n_{P,Q}}"']} \& {}_nQ\otimes{}_nP\handwritten{\arrow[d,"{\alpha^n_{Q,P}}"]}\\
\text{(2.1.9.2)} \& \underline{\mathrm{Tor}}_1(P,Q)\handwritten{\arrow[r,"\sim"]} \& \underline{\mathrm{Tor}}_1(\underset{\deleted{\unsure{n}}}{Q},P)\ ,
\end{tikzcd}}
où les flèches horizontales sont les isomorphismes de symétrie. Ceci posé, les relations entre $\handwritten{\alpha}$ et $\handwritten{\beta}$ sont\handwritten{/}données par la formule:
\[
\handwritten{\beta}^{n}_{P,Q}=-{\handwritten{\alpha}}^{n}_{P,Q}\ ,\tag*{\(\substack{\text{(2.1.9.3)}\\\deleted{\text{\unsure{(2.1}}}}\)}
\]
en d'autres termes:
\margin{\handwritten{\begin{tikzpicture}[baseline]\draw (0,1.1) -- (0,0) -- (.11,0);\end{tikzpicture}}}

\uline{Lemme} 2.1.10. Soient $P$ et $Q$ deux \deleted{\unsure{g|C}}Groupes commutatifs, $n\handwritten{>}0$ un entier. Alors le diagramme suivant est \uline{anticommutatif} , où $\handwritten{\alpha}$ est défini dans (2.1.1)\handwritten{:}
% archive sheet number: 29
